\documentclass{inprogress}
\newcommand{\tr}{\operatorname{tr}}
\newcommand{\sign}{\operatorname{sign}}
\newcommand{\law}{\operatorname{Law}}
\newcommand{\Span}{\operatorname{Span}}
\newcommand{\erfc}{\operatorname{erfc}}


\newlist{hypotheses}{enumerate}{2}
\setlist[hypotheses]{leftmargin = 5em}
\newcommand{\Hyp}[1]{\item[\textsc{#1}]\label{hyp:#1}\!\!\!\textsc{:}}
\newcommand{\hyp}[1]{\textsc{#1}}


\title{A duality approach to $L^2$ stable phase retrieval for independent random variables}

\begin{document}
\maketitle

Let us recall the definition of stable phase retrieval.
\begin{definition}
    Let $E$ be a subspace of $L^p(\mu)$ and fix $C\geq 1$. We say that $E$ does $C$-SPR if for all $f,g\in E$ we have
    \begin{equation}\label{C-SPR}
        \inf_{|\lambda|=1}\|f-\lambda g\|\leq C \| |f|-|g|\|.
    \end{equation}
\end{definition}


In \cite{calderbank2022stable},  Calderbank-Daubechies-Freeman-Freeman were able to prove the following theorem.

\begin{theorem}
 Let $(X_n)_{n=1}^\infty$ be an orthonormal sequence of independent mean-zero random variables in $L^2([0,1])$. The following statements are equivalent.
 \begin{enumerate}
\item The closed span of $(X_n+1_{(n,n+1]})_{n=1}^\infty$ does SPR in $L^2(\mathbb{R}).$
\item The $L^1$ and $L^2$ norms of $(X_n)$ are uniformly comparable, i.e., there exists $A>0$ such that $$A\|X_n\|_{L^2([0,1])}\leq \|X_n\|_{L^1([0,1])}\leq \|X_n\|_{L^2([0,1])}\hspace{3mm} \text{for all}\ n\in \mathbb{N}.$$
\end{enumerate}
\end{theorem}
After writing the paper \cite{calderbank2022stable}, the authors communicated to us the following conjecture.
\begin{conjecture}\label{conjecture CDFF}
    If $(X_n)$ is an orthonormal sequence of mean-zero independent random variables in $L^2$ then span$(X_n)$ does SPR in $L^2$ if and only if there exist constants $0<A\leq B<1$ such that $A\leq \|X_n\|_{L^1} \leq B$ for all but at most one $n\in \mathbb{N}$.
\end{conjecture}
Note that one of the implications of \Cref{conjecture CDFF} is straightforward: If there were no such $A$ then SPR would fail, as it was shown in \cite{FOTP} that if a subspace $E$ of $L^p(\mu)$ ($1\leq p<\infty)$ does SPR then $\|\cdot\|_1\sim \|\cdot\|_p$ on $E$. On the other hand, if such a $B <1$ did not exist then because a mean-zero random variable $f$ has Radamacher distribution if and only if $\|f\|_{L^2}=\|f\|_{L^1}$ we can conclude that $\text{span} (X_n)$
does not perform stable phase retrieval for the same reason that $\text{span} (R_n)$, the span of Rademacher functions, does not.



The purpose of this document is to prove the following theorem.

\begin{theorem}
    \label{thm:goal}
    Let $X_1, \dots, X_n, \dots$ be independent random variables and let $\delta>0$ be such that for all $i\in \mathbb{N}$, we have
    \begin{hypotheses}
        \Hyp{Norm} $\|X_i\|_2 = 1$,
        \Hyp{Mean} $\mathbb E[X_i] =  0$,
        \Hyp{Equiv} $\delta \|X_i\|_2 \le \|X_i\|_1 \le (1-\delta) \|X_i\|_2$.
    \end{hypotheses}
    Then the family of $X_i$ performs SPR in $L^2$ with constant $C=C_\delta$. More specifically, for any $X, Y \in \Span(X_1, \dots, X_n \dots)$ one has \[
    \inf_{|\lambda|=1}\mathbb E[|X-\lambda Y|^2]\le C_\delta \mathbb E[||X|-|Y||^2].
    \]
\end{theorem}


The hypotheses of \Cref{thm:goal} are applied as follows.
\begin{lemma}\label{lem:trunc}
    Let $X$ be a random variable that satisfies the hypotheses \hyp{Norm}, \hyp{Mean} and \hyp{Equiv} in \Cref{thm:goal}. Then there exist $X$-measurable random variables $S,T$ such that

    \begin{minipage}{.49\textwidth}
    \begin{hypotheses}
        \Hyp{ReprS} $\mathbb E[X\cdot S] = 1$,
        \Hyp{MeanS} $\mathbb E[S] = 0$,
        \Hyp{BoundS} $\|S\|_\infty \le C(\delta)$,
        \Hyp{SuppS} $\|X\cdot S\|_\infty \le C(\delta)$,
    \end{hypotheses}
    \end{minipage}
    \begin{minipage}{.49\textwidth}
    \begin{hypotheses}
        \Hyp{ReprT} $\mathbb E[(X^2-1)T] = 1$,
        \Hyp{MeanT} $\mathbb E[T] = 0$,
        \Hyp{BoundT} $\|T\|_\infty \le C(\delta)$,
        \color{red}
        \Hyp{SuppT} $\|X\cdot T\|_\infty \le C(\delta)$.
        \color{black}
    \end{hypotheses}
    \end{minipage}
\end{lemma}
\begin{proof}

By \hyp{Norm}, \hyp{Mean} and \hyp{Equiv} there exists a set $S_+$ of size $\delta$ where $X \in [\delta, \delta^{-1}]$ (and by symmetry, a set $S_{-}$ where $X \in [ -\delta^{-1},-\delta]$). Letting $S:= \frac{\1_{S_+}-\1_{S_-}}{\mathbb E[X\cdot (\1_{S_+}-\1_{S_-})]}$ does the trick. For $T$ just use a balanced variation of $\sign(X^2-1)$ (so that it has mean zero).
\end{proof}
We must also recall the orthogonal reduction from \cite{FOTP}.
\begin{theorem}\label{ortho red}
    Let $f,g\in L^2(\mu)$. Then there exists $f',g'\in \text{span}\{f,g\}$ such that $f'$ is orthogonal to $g'$ and the following two properties hold.
    \begin{enumerate}
        \item The left-hand side of \eqref{C-SPR} is unchanged: $$\inf_{|\lambda|=1}\|f-\lambda g\|=\inf_{|\lambda|=1}\|f'-\lambda g'\|= (\|f'\|^2+\|g'\|^2)^\frac{1}{2}.$$
        \item The right-hand side of \eqref{C-SPR} decreases pointwise:
        $$\left| |f'|-|g'|\right|\leq \left| |f|-|g|\right|.$$
    \end{enumerate}
\end{theorem}
By scaling and invoking \Cref{ortho red}, we see that to prove that a subspace $E$ of $L^2(\mu)$ does SPR it is sufficient to verify that        the inequality \eqref{C-SPR} holds for vectors $f,g\in E$ satisfying $\|f\|=1$, $\|g\|\leq 1$ and $f\perp g$. Note that in this we have $\inf_{|\lambda|=1}\|f-\lambda g\|\in [1,\sqrt{2}]$. This shows that SPR is worst at functions where the quotient norm and the original Hilbert norm agree, and, in particular, we can always assume uniform separation of $f$ and $g$ when proving that a subspace does SPR.





The key estimate is the following representation result:

\begin{lemma}
    Given $\vec a = \{a_k\}_{k=1}^\infty$ with $|\vec a|_2 = 1$, let $X_{\vec a}:= \sum_{j=0}^\infty a_j X_j$ and let $f_{\vec a}:=  \sum_{i<j} S_i S_j a_i a_j + \sum_{j} T_j a_j^2$, where $S_j, T_j$ are defined from $X_j$ as in \Cref{lem:trunc}. Then 
    $$
    |\mathbb E[X_{\vec a}\cdot X_{\vec c}]|^2 = \mathbb E [|X_{\vec a}|^2 \cdot f_{\vec c}].
    $$
\end{lemma}
\begin{proof}
    The independence of the random variables $X_i$ implies that 
    $$|\mathbb E[X_{\vec a}\cdot X_{\vec c}]|^2 = \left|\sum_{i,j} a_i c_j \mathbb E[X_i\cdot X_j]\right|^2  = \left|\sum_{i} a_i c_i\right|^2 =|\vec a \cdot \vec c|^2.$$
    Now combining the independence of the $X_i$ and the mean zero property, we may compute that
    $$
    \begin{aligned}
        \mathbb E [|X_{\vec a}|^2\cdot f_{\vec c}] 
        = &
        \mathbb E [(|X_{\vec a}|^2 -1)\cdot f_{\vec c}] 
        \\= &
        \sum_{i\neq j} \mathbb E[X_{i}X_j S_{ij}] a_i a_j c_i c_j+
        \sum_{j} \mathbb E[(X_j^2-1) T_j] a_j^2 c_j^2
        \\= &
        \sum_{i\neq j} a_i a_j c_i c_j+
        \sum_{j}  a_j^2 c_j^2
        \\= &
        |\vec a \cdot \vec c|^2.
    \end{aligned}
    $$
\end{proof}

If the random variables $X_i$ were uniformly in $L^{2+2\epsilon}$ and $X, Y$ are $L^2$-normalized such that $\|X-Y\|_{2}\gtrsim 1$, one could finish directly by using that 
$$
\begin{aligned}
    2 = \|X-Y\|_2^2 \approx \max_{\|\vec c\|_{l^2 = 1}} |\mathbb E[X X_{\vec c}]|^2- |\mathbb E[Y X_{\vec c}]|^2
    = 
    \max_{\|\vec c\|_{l^2 = 1}} |\mathbb E[(|X|^2-|Y|^2) f_{\vec c}]|.
\end{aligned}
$$
Indeed, applying Holder's inequality
where $\frac 1  M + \frac{1}{1+\epsilon} = 1$ and $ |\mathbb E[(|X|^2-|Y|^2) ]| \le \||X|-|Y|\|_{2} \||X|+|Y|\|_{2}\le 2  \||X|-|Y|\|_{2}$ one obtains that
$$
    1\lesssim |\mathbb E[(|X|-|Y|) ]| \|f_{\vec c}\|_M,
$$
where $\|f_{\vec c}\|_M$ is bounded by sub-Gaussian bounds on sums of $L^\infty$ independent mean-zero random variables. The main difficulty remaining in the $L^2$ case is solved by a summation by parts argument. 


\begin{proof}[Proof of \Cref{thm:goal}]




Let $\phi$ be a smooth cut-off function such that $\phi([-1/2,1/2]) = 0$ and $\phi([-1,1]^c) = 1$. Then
$$
\begin{aligned}
        ||\mathbb E[X_{\vec a} X_{\vec c}]|^2- |\mathbb E[X_{\vec b} X_{\vec c}]|^2| 
        =& 
        |\mathbb E[(|X_{\vec a}|^2-|X_{\vec b}|^2) f_{\vec c}]|
        \\=&
        |\mathbb E[(|X_{\vec a}|^2-|Y_{\vec b}|^2) f_{\vec c}(\phi(M^{-1}f\vec c) + 1 - \phi(M^{-1}f_{\vec c}))]|
        \\\le&
        M \mathbb E[||X_{\vec a}|^2-|X_{\vec b}|^2|] 
        + |\mathbb E[(|X_{\vec a}|^2-1) f_{\vec c}\phi(M^{-1} f_c)]|
        + |\mathbb E[(|X_{\vec b}|^2-1) f_{\vec c}\phi(M^{-1} f_c)]|.
\end{aligned}
$$
The remaining goal is to bound the second and third terms by an arbitrarily small constant when $M$ is large enough, so that they can be absorbed into the left-hand side. Since both terms are equivalent, we focus on bounding the second one. Let $S_{ij}:= S_i S_j$ when $i\neq j$ and $S_{ii} = T_i$. Similarly, let $X_{ij}= X_i X_j$ when $i\neq j$, and $X_{ii} = X_{i}^2 - 1$. Then, if $|\vec a|^2=1$, we have

$$
X_{\vec a}^2 -1 = \sum_{i,j} a_i a_j X_{ij} \quad \text{and} \quad f_{\vec c} = \sum_{k,l} c_i c_j S_{kl}.
$$

We will split the sum into several terms.

\begin{enumerate}
\item\textbf{Sums when $k\neq l$.}
\begin{enumerate}
    \item In this first case, assume that neither $i$ nor $j$ are equal to any of $k,l$. Define $X_{\hat k,\hat l}:= \sum_{\substack{i\neq k,l\\j\neq k,l}} a_{i} a_jX_{ij}$. The variable $X_{\hat k,\hat l}$ has mean zero, is independent of $X_{k,l}$, and satisfies $\|X_{\hat k,\hat l}\|_{L^1}\lesssim 1$. From here, by summation by parts one deduces
    $$| \mathbb E[X_{\hat k\hat l}S_{kl}\phi(M^{-1}f_{\vec c})]| \lesssim M^{-2}|c_k c_l|.$$
    This is because one has a uniform bound $\mathbb E_{k,l}[S_{kl}\phi(M^{-1}f_{\vec c})]\lesssim \|\partial_{c_i}\partial_{c_j} \phi(M^{-1}f_{\vec c})\|_\infty \lesssim M^{-1} c_i c_j$. A more careful estimate should give us an extra exponential gain in $M$ (which should give us better dependency on $\delta$) but for a qualitative result this works. This, in particular, shows that 
    $$
    \left|\sum_{k}\sum_{l\neq k} \sum_{i\neq k,l}\sum_{j\neq k,l} a_i a_j c_k c_l \mathbb E[X_{ij}S_{kl}\phi(M^{-1}f_{\vec c})] \right| \lesssim M^{-2}.
    $$
    \item We now focus on the case when $i=k, j=l$ (or flipped, by adding a factor of two). In this case, we have
    $$
    |\mathbb E[S_i X_i S_j X_j\phi(M^{-1}f_{\vec c})]|
    \lesssim 
    \|S_i X_i\|_\infty \|S_j X_j\|_\infty e^{-M^2/2},
    $$ 
    and then we can sum.
    \item Finally, we are left with the intermediate case where $i=k$ and $j\neq l$. We look at the generic case $i=k\neq l \neq j$. We have a sum
    $$
    a_i a_j c_j c_k \mathbb E[X_i X_j S_j S_k]
    $$
    and we take the expected value over $i,k$ to witness a $M^{-2}c_ic_k$ gain. 
\end{enumerate}
\item \textbf{Sums when $k=l$.}
\begin{enumerate}
    \item Now, if instead neither $i,j$ are equal to $k=l$, one uses summation by parts to obtain the pointwise bound 
    $$
     | \mathbb E_k [X_{\hat k\hat k}S_{kk}\phi(M^{-1}f_{\vec c})]| \lesssim M^{-1}|c_k \cdot X_{\hat k\hat k}| \|S_{kk}\|_2,
    $$
    and therefore by Jensen
    $$| \mathbb E[X_{\hat k\hat k}S_{kk}\phi(M^{-1}f_{\vec c})]| \lesssim M^{-1}|c_k|.$$
 This gives us the bound 
    $$
    \left|\sum_{k}\ \sum_{i\neq k}\sum_{j\neq k} a_i a_j c_k^2 \mathbb E[X_{ij}S_{kk}\phi(M^{-1}f_{\vec c})] \right|
    =
    \left|\sum_{k} c_k^2 \mathbb E[X_{\hat k \hat k}S_{kk}\phi(M^{-1}f_{\vec c})] \right|
    \lesssim M^{-1}.
    $$
    \item  In the case $i \neq j = k = l$,  we have
    $$
    a_i a_j c_j^2 \mathbb E[X_{i} X_j S_{jj}\phi(M^{-1}f_{\vec c})],
    $$
    and taking the expected value over $i$ we gain a $M^{-1}c_i$:
    $$
    |\mathbb E_i[X_{i} X_j S_{jj}\phi(M^{-1}f_{\vec c})]| \lesssim |c_i M_i^{-1} X_j|  \|X_i\|_2
    $$
    
    and we can sum.
    \item There is the last extreme case $i=j=k=l$. In this case we have
    $$
    \sum_j a_j^2 c_j^2 \mathbb E[|X_j^2 -1| \phi(M^{-1}f_{\vec c}) ]
    $$
    this term we actually treat as a difference instead
    $$
    \sum_j (a_j^2-b_j^2) c_j^2 \mathbb E[|X_j^2 -1| \phi(M^{-1}f_{\vec c}) ]
    \lesssim\max_j |a_j^2-b_j^2| = \max_j \mathbb E[T_{jj}||X_{\vec a}|^2-|X_{\vec b}|^2|]
    $$

    
\end{enumerate}
\end{enumerate}



\begin{lemma}[Almost-independent random variables]
    Let $X$ be a mean-zero random variable, and $g$ be a differentiable function. Then 
    $$
    \mathbb E[X g(X)] \le \|g'\|_{\infty}\|X\|_2.
    $$
\end{lemma}
\begin{proof}
    WLOG $g(0)=0$. The proof is then concluded by saying indeed.
\end{proof}



\end{proof}

\begin{conjecture}
Let $X$ be a CLT-limit of the sequence of random variables $(X_{n,k})_{1\le k \le n <\infty}$. Then $X$ is a CLT-limit of (normalized) Gaussians and Poissons. 
\end{conjecture}


\bibliographystyle{plain}
\bibliography{refs.bib}

\end{document}
